% Figure: Jeffcott rotor whirl amplitude and phase through the first
% critical speed. Amplitude (normalised to eccentricity) rises to a peak
% near r=1 and settles to unity above; phase runs 0 -> 90 -> 180 deg.
% Two y-axes emulated by scaling the phase curve into the same box and
% labelling it.
\begin{figure}[htbp]
\centering
\begin{tikzpicture}
\begin{axis}[
  width=12cm, height=7cm,
  xlabel={speed ratio $r=\Omega/\omega_n$},
  ylabel={whirl amplitude $X/e$},
  xmin=0, xmax=3, ymin=0, ymax=4.2,
  axis lines=left,
  domain=0.001:3, samples=400,
  legend pos=north east,
  legend cell align=left,
  every axis x label/.style={at={(current axis.right of origin)}, anchor=west},
  every axis y label/.style={at={(current axis.above origin)}, anchor=south},
  grid=major, grid style={enggray!20},
  tick label style={font=\scriptsize},
  label style={font=\small},
  legend style={font=\scriptsize},
]
% Jeffcott whirl amplitude: X/e = r^2 / sqrt((1-r^2)^2 + (2 zeta r)^2)
\addplot[engblue, very thick] {x^2/sqrt((1-x^2)^2 + (2*0.08*x)^2)};
\addlegendentry{$X/e$, $\zeta=0.08$}
% high-speed self-balancing asymptote at 1
\addplot[enggray, thin, forget plot] {1};
\node[enggray, font=\scriptsize, anchor=south west] at (axis cs:2.05,1.05)
  {self-balancing limit $X/e\to 1$};
% phase curve scaled: phase/180*4 so 180deg maps to 4 on the axis.
% phase = atan2(2 zeta r, 1 - r^2) in degrees, /180*4.
\addplot[engred, thick, dashed]
  {(atan2(2*0.08*x, 1 - x^2))/180*4};
\addlegendentry{phase lag (right scale)}
% right-scale tick labels for the phase curve
\node[engred, font=\scriptsize, anchor=east] at (axis cs:2.98,0.05) {$0^{\circ}$};
\node[engred, font=\scriptsize, anchor=east] at (axis cs:2.98,2.05) {$90^{\circ}$};
\node[engred, font=\scriptsize, anchor=east] at (axis cs:2.98,4.05) {$180^{\circ}$};
\end{axis}
\end{tikzpicture}
\caption[Whirl amplitude and phase of a Jeffcott rotor]{The whirl
response of a Jeffcott rotor, a disc with mass-centre eccentricity $e$
on a flexible shaft, driven by its own rotating unbalance as the speed
$\Omega$ sweeps through the first bending critical $\omega_n$. The whirl
amplitude $X/e$, blue, follows the rotating-unbalance law $r^2/\sqrt{(1
-r^2)^2+(2\zeta r)^2}$: near zero at low speed, a resonant peak at the
critical, and a settling to unity above it. Above the critical the disc
whirls about its own mass centre rather than its geometric centre, the
self-balancing that keeps a supercritical machine quiet once it is
through the resonance. The phase lag, red dashed against the right
scale, runs from $0^{\circ}$ below the critical through $90^{\circ}$ at
$r=1$ to $180^{\circ}$ above, the same phase march as the forced
oscillator and the signature that locates the critical in a run-up
test.}
\label{fig:vol03ch06:whirl-orbit}
\end{figure}
